\documentclass[11pt]{amsart}
\usepackage{geometry}
\geometry{a4paper}
\usepackage{amssymb}
\usepackage[dvipsnames]{xcolor}
\usepackage[colorinlistoftodos,prependcaption]{todonotes}

\newtheorem{theorem}{Theorem}
\newtheorem{lemma}[theorem]{Lemma}
\newtheorem{proposition}[theorem]{Proposition}
\newtheorem{conjecture}[theorem]{Conjecture}

\theoremstyle{definition}
\newtheorem{definition}{Definition}
\newtheorem{example}{Example}

\theoremstyle{remark}
\newtheorem{remark}[theorem]{Remark}

% ── Macros shared with agg-logic.tex ─────────────────────────────────────────
\newcommand{\DataTypes}{\mathcal{T}}
\newcommand{\DataValues}{\mathcal{V}}
\newcommand{\DomObjects}{\mathcal{O}}
\newcommand{\DomValues}{\mathcal{O}^{\mathsf{val}}}
\newcommand{\DomEnts}{\mathcal{O}^{\mathsf{ent}}}
\newcommand{\Preds}{\mathcal{P}}
\newcommand{\EntPreds}{\Preds^{\textsf{ent}}}
\newcommand{\ValPreds}{\Preds^{\textsf{val}}}
\newcommand{\FactPreds}{\Preds^{\textsf{fact}}}
\newcommand{\RefSchemes}{\mathcal{R}}
\newcommand{\sem}[1]{[\![{#1}]\!]}
\newcommand{\pr}[1]{\mathsf{\textcolor{RoyalBlue}{#1}}}
\newcommand{\var}[1]{\mathsf{\textcolor{Plum}{#1}}}
\newcommand{\tp}[1]{\texttt{\textsc{#1}}}
\newcommand{\dataOp}[1]{\mathbin{\textcolor{ForestGreen}{#1}}}
\newcommand{\dataVal}[1]{\textcolor{ForestGreen}{#1}}
\newcommand{\dataEq}{\mathrel{\textcolor{ForestGreen}{\doteq}}}

% ── New in this document ─────────────────────────────────────────────────────
% Frontier (data) variables, shared between the two sides of a mapping rule.
\newcommand{\fv}[1]{\mathit{\textcolor{BrickRed}{#1}}}
\newcommand{\NULL}{\dataVal{\mathsf{null}}}
\newcommand{\DataValuesN}{\DataValues_{\mathsf{null}}}
\newcommand{\Opt}[1]{\mathsf{opt}_{#1}}
\newcommand{\IdSchemes}{\mathcal{I}}
\newcommand{\ident}{\mathit{id}}
\newcommand{\corr}{\;\Longleftrightarrow\;}
\newcommand{\Cov}{\mathsf{Cov}}
\newcommand{\Det}{\mathsf{Det}}

\title{Schema mappings for FreeFBM: tree-pattern correspondences}
\author{Jan Hidders}
\date{November 2026 (draft)}

\begin{document}
\maketitle

\setcounter{tocdepth}{1}
\tableofcontents

\section{Goal}

FreeFBM is to become a workbench for designing, redesigning and reverse
engineering database schemas, with ORM2 as the common language. A
relational schema is represented as an ORM2 schema of a restricted form (a
\emph{profile}), and the relationship between two schemas is documented by
a \emph{mapping}: a finite set of rules, each relating a pattern over one
schema to a pattern over the other.

This document defines a candidate language for these rules. It is chosen to
be expressive enough for the usual relational design transformations
(vertical and horizontal partitioning, foreign keys versus cross tables,
nullable columns) while keeping three properties:
\begin{enumerate}
  \item \emph{No value invention.} Every object on one side that does not
    occur on the other is determined by key and identification constraints
    from values that do.
  \item \emph{Symmetry.} A rule is a correspondence, not a directed
    implication, so the inverse of a mapping is obtained by swapping the two
    sides of every rule.
  \item \emph{Checkability.} Whether a mapping is a bijection between the
    legal instances of the two schemas reduces to a small number of
    constraint implications (key-safety, coverage), and the constraints
    that would have to be added to either schema to make it one can be
    reported to the user (\emph{constraint accounting}).
\end{enumerate}
The logic is that of \emph{An intermediate sorted logic for ORM2}
(\texttt{agg-logic.tex}); we reuse its notation and only add what mappings
need. Aggregation, transitive closure and data functions are deliberately
left out of the mapping language: a rule that needs them is expressed by a
derived fact type inside one of the schemas, which the mapping then relates
with a simple rule.

\section{Setting}
\label{sec:setting}

\subsection{Schemas}

A \emph{schema} $S = (V_S, \Sigma_S)$ consists of a vocabulary $V_S =
(\Preds_S, \delta_S, \sigma_S, \RefSchemes_S)$ as in \texttt{agg-logic.tex}
and a set $\Sigma_S$ of constraints, which are closed formulas of that
logic. An \emph{instance} of $S$ is an interpretation $\iota_S$ of $V_S$
with $\iota_S \models \Sigma_S$. A mapping relates two schemas $S$ and $T$
whose predicate sets are disjoint, so that an instance of $S$ and an
instance of $T$ together form an interpretation $\iota_S \cup \iota_T$ of
the combined vocabulary.

\subsection{Identification schemes}

The vocabulary of \texttt{agg-logic.tex} only has simple reference schemes
$\RefSchemes(p) = (r, q)$: an entity is identified by one value through one
binary fact type. Relational keys and ORM2 external uniqueness constraints
need composite identification, so we generalise.

\begin{definition}
An \emph{identification scheme} for an entity type $p \in \EntPreds$ is a
non-empty sequence $(r_1, \ldots, r_k)$ of binary fact types with
$\sigma(r_i) = (p, q_i)$ and $q_i \in \EntPreds \cup \ValPreds$. A schema
\emph{uses} it if $\Sigma$ implies that in every instance the relation
$$\{ (e, o_1, \ldots, o_k) \mid e \in \iota(p),\ (e, o_i) \in \iota(r_i)
  \text{ for } 1 \le i \le k \}$$
is a total injective function from $\iota(p)$ to $\iota(q_1) \times \cdots
\times \iota(q_k)$. That is: each $r_i$ is mandatory and functional for
$p$, and the $r_i$ together are unique (an external uniqueness constraint,
the preferred identifier of $p$). We write $\IdSchemes(p)$ for the
identification scheme used for $p$, if any, and require the graph with an
edge from $p$ to each $q_i$ to be acyclic.
\end{definition}

A simple reference scheme $\RefSchemes(p) = (r, q)$ is the special case
$\IdSchemes(p) = (r)$ with $q \in \ValPreds$. By acyclicity, every entity
$e$ of a type with an identification scheme has an \emph{identifying tuple}
$\ident(e) \in \DataValues^{+}$ of data values, obtained by recursively
replacing every identifying entity by its own identifying tuple and every
domain value $o$ by $\rho(o)$. Two entities of the same type are equal if
and only if their identifying tuples are equal.

\begin{remark}
ORM2 requires a preferred identifier for every entity type, so in a
validated FreeFBM schema every entity type has an identification scheme,
possibly inherited from a supertype. The role order ``entity first'' in
$\sigma(r_i)$ is for presentation only; FreeFBM allows the identifying role
in any position.
\end{remark}

\subsection{The relational profile}

A relational schema is represented in \emph{encoding~B}: every table is an
entity type, every column is a binary fact type from that entity type to a
value type (or, for an abstract identifier column, to an entity type), a
\texttt{NOT NULL} column is a mandatory role, the primary key is the
identification scheme of the table's entity type, a \texttt{UNIQUE}
constraint is an external uniqueness constraint, and a foreign key is a
subset constraint between role sequences. Relational notation is an
abbreviation over this encoding (Definition~\ref{def:relatom}).

\subsection{Frontier variables and $\NULL$}

Besides the object variables of \texttt{agg-logic.tex}, which range over
the domain objects of a type, mapping rules use \emph{frontier variables},
written $\fv{y}, \fv{y}_1, \ldots$, which range over $\DataValuesN =
\DataValues \cup \{ \NULL \}$, where $\NULL \notin \DataValues$. Frontier
variables are the only variables shared by the two sides of a rule, so all
that passes between two schemas is data values. Within mapping rules
$\NULL$ is an ordinary constant: $\NULL \dataEq \NULL$ holds (as SQL's
\texttt{IS NOT DISTINCT FROM}), and the logic stays two-valued.

\section{Patterns}
\label{sec:patterns}

\begin{definition}
\label{def:pattern}
A \emph{pattern} over a vocabulary $V$ is a formula built by the grammar
$$\varphi \;::=\; p(u_1, \ldots, u_n) \;\mid\; q(u) \;\mid\; u \dataEq \fv{y}
  \;\mid\; u \dataEq c \;\mid\; \varphi \wedge \varphi \;\mid\;
  \Opt{\fv{y}}(\varphi) \;\mid\; \neg \exists \bar u\, \varphi$$
where $p \in \FactPreds$, $q \in \EntPreds \cup \ValPreds$, the $u_i$ are
object variables, $\fv{y}$ is a frontier variable and $c$ a data constant,
subject to the following conditions.
\begin{enumerate}
  \item[(P1)] \emph{Acyclicity.} The hypergraph whose nodes are the object
    variables of $\varphi$ and whose hyperedges are the variable sets of its
    fact atoms is $\alpha$-acyclic, both at the top level and within every
    $\Opt{}$ and $\neg\exists$ part together with the variables it shares
    with its context.
  \item[(P2)] \emph{Binding of frontier variables.} A binding atom
    $u \dataEq \fv{y}$ has a value-type variable $u$, or an entity-type
    variable whose type has a simple reference scheme. Every frontier
    variable occurring in $\varphi$ occurs in a binding atom at the top
    level, or is the \emph{guard} $\fv{y}$ of exactly one optional part
    $\Opt{\fv{y}}(\varphi')$. A guard occurs in a binding atom inside its
    own optional part and nowhere outside it. A frontier variable that is
    not a guard may occur in several binding atoms, which then join on its
    value.
  \item[(P3)] \emph{Guarded negation.} In $\neg\exists \bar u\, \varphi'$,
    every free variable of $\varphi'$ not in $\bar u$ is an object variable
    occurring in a fact or type atom of the enclosing pattern outside the
    negation, or a frontier variable bound at the top level.
\end{enumerate}
The object variables occurring at the top level of $\varphi$, outside every
$\Opt{}$ and $\neg\exists$ part, are its \emph{local variables}. Variables
introduced inside an $\Opt{}$ part are local to that part and existentially
quantified there.
\end{definition}

\begin{definition}
The optional part abbreviates a first-order formula. Let $\bar w$ be the
object variables that occur in $\varphi'$ but not in its context. Then
$$\Opt{\fv{y}}(\varphi') \;:=\;
  \big(\fv{y} \dataEq \NULL \,\wedge\, \neg \exists \bar w\, \exists \fv{y}'\,
     \varphi'[\fv{y}'/\fv{y}]\big)
  \;\vee\;
  \big(\neg(\fv{y} \dataEq \NULL) \,\wedge\, \exists \bar w\, \varphi'\big).$$
\end{definition}

The first disjunct says that the optional branch is absent and the value is
$\NULL$; the second that it is present with value $\fv{y}$. The negation
$\neg\exists \bar u\,\varphi'$ is the general absence test (an anti-join),
and is safe by (P3).

\begin{definition}
\label{def:relatom}
Let $R$ be a table entity type with columns $\mathit{col}_1, \ldots,
\mathit{col}_n$ (binary fact types with $\sigma(\mathit{col}_i) = (R,
q_i)$). A \emph{relational atom} $R(t_1, \ldots, t_n)$, where each $t_i$ is
a frontier variable $\fv{y}$, a nullable frontier variable $\fv{y}?$, a
constant $c$, the absence marker $-$, the wildcard $\_$, or an object
variable $u$ (for a column whose $q_i$ is an entity type), abbreviates
$$\exists x \big( R(x) \wedge \textstyle\bigwedge_{i=1}^{n} \alpha_i(x)
  \big)$$
with a fresh local variable $x$ of type $R$ and fresh variables $v_i$ of
type $q_i$, where
\begin{center}
\begin{tabular}{ll}
$t_i = \fv{y}$: & $\alpha_i = \exists v_i\,(\mathit{col}_i(x, v_i) \wedge
  v_i \dataEq \fv{y})$\\
$t_i = \fv{y}?$: & $\alpha_i = \Opt{\fv{y}}(\mathit{col}_i(x, v_i) \wedge
  v_i \dataEq \fv{y})$\\
$t_i = c$: & $\alpha_i = \exists v_i\,(\mathit{col}_i(x, v_i) \wedge
  v_i \dataEq c)$\\
$t_i = -$: & $\alpha_i = \neg\exists v_i\, \mathit{col}_i(x, v_i)$\\
$t_i = \_$: & $\alpha_i = \top$\\
$t_i = u$: & $\alpha_i = \mathit{col}_i(x, u)$.
\end{tabular}
\end{center}
A negated relational atom $\neg R(t_1, \ldots, t_n)$ abbreviates the
negation of this formula, with the fresh variables quantified inside the
negation.
\end{definition}

\begin{remark}
\label{rem:optrow}
Two different kinds of optionality occur in relational examples. A
\emph{nullable column} $R(\fv{a}, \fv{b}, \fv{c}?)$ requires the row to
exist and makes only the column optional. An \emph{optional row}
$\Opt{\fv{c}}(R_3(d, \fv{c}))$ makes the whole atom optional: $\fv{c}$ is
$\NULL$ if and only if no row of $R_3$ with first column $d$ exists. In
encoding~B, where a row is an entity and each column a role, the first is
an optional role of an existing entity and the second an optional entity.
Earlier informal notes wrote both as $R(\ldots, \fv{c}?)$; the examples
below use $\Opt{}$ for the second.
\end{remark}

The local variable $x$ is the row. For rules over relational schemas it is
always determined by the primary key (see Section~\ref{sec:keysafety}), so
it never needs to be invented.

\section{Rules and mappings}
\label{sec:rules}

\begin{definition}
A \emph{rule} from $S$ to $T$ has the form
$$\varphi(\bar x, \bar{\fv{y}}) \corr \psi(\bar{\fv{y}}, \bar z)$$
where $\varphi$ is a pattern over $V_S$ with local variables $\bar x$,
$\psi$ is a pattern over $V_T$ with local variables $\bar z$, and
$\bar{\fv{y}}$ is the set of frontier variables, each of which occurs in
both patterns. Its meaning is the closed formula
$$\sem{\varphi \corr \psi} \;:=\; \forall \bar{\fv{y}} \,\big[\,
  \exists \bar x\, \varphi \;\leftrightarrow\; \exists \bar z\, \psi
  \,\big],$$
where $\forall \bar{\fv{y}}$ ranges over $\DataValuesN$. A \emph{mapping}
$M = (S, T, \Gamma)$ consists of two schemas with disjoint vocabularies and
a finite set $\Gamma$ of rules from $S$ to $T$. It denotes the relation
$$\sem{M} \;=\; \{ (\iota_S, \iota_T) \mid \iota_S \models \Sigma_S,\
  \iota_T \models \Sigma_T,\ \iota_S \cup \iota_T \models
  \textstyle\bigwedge_{\gamma \in \Gamma} \sem{\gamma} \}.$$
\end{definition}

\begin{lemma}[Finite range]
For a rule $\varphi \corr \psi$ and instances $\iota_S, \iota_T$, if
$\exists \bar x\,\varphi$ holds for an assignment to $\bar{\fv{y}}$, then
every $\fv{y}$ is assigned $\NULL$ or the data value of a domain object in
$\iota_S$; likewise for $\psi$ and $\iota_T$. Hence $\sem{\varphi \corr
\psi}$ can be decided by evaluating both patterns over finitely many
assignments.
\end{lemma}

\begin{proof}
By (P2), each $\fv{y}$ occurs in a positive binding atom $u \dataEq
\fv{y}$ at the top level, or is the guard of an optional part. In the first
case, and in the present case of the optional part, $u$ ranges over the
domain objects of $\iota_S$. In the absent case the value is $\NULL$.
\end{proof}

\begin{proposition}[Inverse]
Let $M^{-1} = (T, S, \{ \psi \corr \varphi \mid \varphi \corr \psi \in
\Gamma \})$. Then $\sem{M^{-1}} = \{ (\iota_T, \iota_S) \mid (\iota_S,
\iota_T) \in \sem{M} \}$.
\end{proposition}

\begin{proof}
Immediate, since $\leftrightarrow$ is symmetric.
\end{proof}

A mapping is meant to relate instances only up to the choice of domain
entities, which carry no information of their own. We write $\iota \cong
\iota'$ if there is a bijection between their domain objects that preserves
types, facts and the data values $\rho$ of domain values. Since the rules
only exchange data values, $\sem{M}$ is closed under $\cong$ on each side.
The question that matters is therefore whether $\sem{M}$ is a bijection
\emph{up to} $\cong$.

\section{Key-safety}
\label{sec:keysafety}

\begin{definition}
Let $\varphi$ be a pattern over a schema $(V, \Sigma)$ with frontier
variables $\bar{\fv{y}}$. The set $\Det(\varphi)$ of \emph{determined}
local variables is the smallest set $D$ such that:
\begin{enumerate}
  \item[(D1)] if $u \dataEq \fv{y}$ or $u \dataEq c$ occurs at the top level
    of $\varphi$, then $u \in D$;
  \item[(D2)] if $e$ is a local variable of an entity type $p$ with
    $\IdSchemes(p) = (r_1, \ldots, r_k)$, and $\varphi$ contains atoms
    $r_i(e, u_i)$ with $u_i \in D$ for all $i$, then $e \in D$;
  \item[(D3)] if $\Sigma$ contains a uniqueness constraint over the
    positions $K$ of fact type $p$, and $\varphi$ contains an atom $p(u_1,
    \ldots, u_n)$ with $u_k \in D$ for all $k \in K$, then $u_j \in D$ for
    all $j$;
  \item[(D4)] more generally, if $\Sigma$ implies a \emph{pattern key}
    $\forall \bar v\, \forall w, w'\, (\pi(\bar v, w) \wedge \pi(\bar v, w')
    \to w = w')$ for a conjunctive pattern $\pi$, and $\varphi$ contains an
    instance of $\pi$ whose $\bar v$ are in $D$, then its $w$ is in $D$.
\end{enumerate}
The pattern is \emph{key-safe} if every local variable is in
$\Det(\varphi)$, and, for every optional part $\Opt{\fv{y}}(\varphi')$,
every variable local to $\varphi'$ is determined in $\varphi'$ starting
from $D \cup \{ \fv{y} \}$. Variables under a negation need not be
determined. A rule is key-safe if $\varphi$ is key-safe over $S$ and $\psi$
over $T$; a mapping is key-safe if all its rules are.
\end{definition}

(D2) and (D3) are special cases of (D4); they are listed because they are
what a checker in FreeFBM would mostly use: preferred identifiers, internal
uniqueness constraints and external uniqueness constraints.

\begin{lemma}[Unique witness]
If $\varphi$ is key-safe and $\iota \models \Sigma$, then for every
assignment to $\bar{\fv{y}}$ there is at most one assignment to the local
variables $\bar x$ that satisfies the top level of $\varphi$.
\end{lemma}

\begin{proof}[Proof sketch]
By induction on the derivation of $u \in \Det(\varphi)$. (D1) uses that
$\rho$ is injective on each value type (and on reference values for entity
types with a simple reference scheme), (D2) the injectivity of
identification, (D3) and (D4) the key constraints. Optional parts are
handled the same way relative to their context.
\end{proof}

So under key-safety, $\exists \bar x\,\varphi$ is in effect a definite
description: ``the'' $\bar x$ with $\varphi$. This is what lets the
declarative reading also be used to \emph{build} the other side: for each
assignment to $\bar{\fv{y}}$ satisfying $\varphi$, the witness of $\psi$
is constructed, and each new entity is the one with the identifying tuple
that the binding gives it.

\section{Coverage and bijectivity}
\label{sec:coverage}

The rules say nothing about a fact that is part of no match. For instance,
in Example~\ref{ex:vertical} a row of $R_2$ without a row of $R_1$ for the
same $d$ completes no match, and the conjunction of rules holds whether or
not that row is there.

\begin{definition}
Let $\psi$ occur in a rule. An atom occurrence $p(u_1, \ldots, u_n)$ or
$q(u)$ at the top level of $\psi$, or inside an optional part, \emph{covers}
the fact $(p, o_1, \ldots, o_n)$ of an instance $\iota$ if some satisfying
assignment of $\psi$ in $\iota$ maps each $u_i$ to $o_i$. A mapping is
\emph{$T$-covering} if in every instance $\iota_T$ of $T$, every fact and
every object of every type is covered by some atom occurrence in some rule.
\end{definition}

$T$-covering is the implication $\Sigma_T \models \Cov_T$, where $\Cov_T$
is the conjunction, over all predicates $p$ of $V_T$, of
$$\forall \bar u\, \big( p(\bar u) \to \textstyle\bigvee_{(\gamma, a)}
  \exists \bar{\fv{y}}\, \exists (\bar z \setminus \bar u)\, \psi_\gamma
  [a \mapsto \bar u] \big),$$
with the disjunction over all rules $\gamma$ and occurrences $a$ of $p$ in
the target pattern $\psi_\gamma$. $S$-covering is defined symmetrically.

\begin{conjecture}[Bijection up to isomorphism]
\label{conj:bijection}
Let $M$ be key-safe, $S$-covering and $T$-covering, and let every entity
type occurring in a rule have an identification scheme. If $(\iota_S,
\iota_T), (\iota_S', \iota_T') \in \sem{M}$, then $\iota_S \cong \iota_S'$
if and only if $\iota_T \cong \iota_T'$.
\end{conjecture}

\begin{proof}[Proof idea]
For each rule $\gamma$ let $Y_\gamma(\iota_S)$ be the set of assignments to
$\bar{\fv{y}}$ satisfying $\exists \bar x\,\varphi_\gamma$. By the
semantics, $Y_\gamma(\iota_S) = Y_\gamma(\iota_T)$. By coverage every fact
of $\iota_S$ lies in a match, and by unique witness each match is
determined by its assignment. Identification makes the choice of entities
consistent across rules: an entity occurring in two matches has the same
identifying tuple in both. So $\iota_S$ is determined up to $\cong$ by the
family $(Y_\gamma(\iota_S))_\gamma$, and likewise $\iota_T$.
\todo{Check the identification argument when an entity is determined only
by (D3) or (D4) and not by its own identification scheme.}
\end{proof}

The conjecture gives injectivity. Totality, meaning that every instance of
$S$ has a corresponding instance of $T$, additionally requires that the
instance built from the rules satisfies $\Sigma_T$. That is exactly what
constraint accounting checks.

\section{Constraint accounting}
\label{sec:accounting}

For a mapping satisfying the conditions of
Conjecture~\ref{conj:bijection}, each predicate of $S$ is definable over
$T$, with entities replaced by their identifying tuples: an $S$-fact holds
if and only if some rule match produces it, and the match exists on the $S$
side if and only if its target pattern holds on the $T$ side. Write
$\chi^{\to T}$ for the result of rewriting a constraint $\chi$ over $S$ in
this way, by unfolding every atom through the rules that produce it.
Every constraint of $\Sigma_S$ then gets one of four statuses with respect
to $T$:
\begin{description}
  \item[native] $\chi^{\to T}$ is (equivalent to) a constraint in
    $\Sigma_T$, for a relational target typically a \texttt{PRIMARY KEY},
    \texttt{UNIQUE}, \texttt{FOREIGN KEY}, \texttt{NOT NULL} or
    \texttt{CHECK};
  \item[implied] $\Sigma_T \models \chi^{\to T}$, shown syntactically;
  \item[lost] a counterexample exists: an instance of $T$ violating
    $\chi^{\to T}$; the user may add $\chi^{\to T}$ to $T$ as a formula
    constraint or assertion, or mark it deontic;
  \item[unknown] neither a proof nor a counterexample was found.
\end{description}
The same is done for $\Sigma_T$ with respect to $S$, and for the coverage
conditions $\Cov_S$ and $\Cov_T$, whose missing parts are reported in the
same way: these are the constraints that ``should have been added if the
schemas are to be equivalent''.

\begin{remark}
Implication between ORM2 constraints is undecidable in general, so the
statuses are not complete. For keys, mandatory roles, subset and exclusion
constraints over tree patterns a decision procedure is plausible, and in
the examples below everything is decided by unfolding plus simple key and
inclusion reasoning.
\end{remark}

\section{Examples}
\label{sec:examples}

The examples use relational notation (Definition~\ref{def:relatom}).
Columns marked $\#$ hold abstract identifiers: they are entity-valued, so
in a rule they appear as object variables, never as frontier variables.

\begin{example}[Vertical partitioning]
\label{ex:vertical}
$S_1$ has $R(a, b, c?)$ with key $\{a, b\}$. $S_2$ has $R_1(d\#, a)$,
$R_2(d\#, b)$ and $R_3(d\#, c)$, each with key $\{d\}$. The mapping has one
rule:
$$R(\fv{a}, \fv{b}, \fv{c}?) \corr R_1(d, \fv{a}) \wedge R_2(d, \fv{b})
  \wedge \Opt{\fv{c}}(R_3(d, \fv{c})).$$
On the left, $\fv{c}?$ is a nullable column; on the right, the row of $R_3$
is optional (Remark~\ref{rem:optrow}), and $d$ is a local variable. Key-safety on the right needs $d$ to be
determined by $\fv{a}$ and $\fv{b}$. The keys $\{d\}$ of $R_1$ and $R_2$ go
the other way, so $S_2$ needs the pattern key ``$\{a, b\}$ determines $d$
in $R_1(d, a) \wedge R_2(d, b)$'', which in ORM2 is an external uniqueness
constraint and the preferred identifier of the type of $d$. $S_2$-covering
needs every $d$ in $R_1$ to also be in $R_2$ and vice versa, and every $d$
in $R_3$ to be in $R_1$. With these three constraints added to $S_2$, the
mapping satisfies the conditions of Conjecture~\ref{conj:bijection}.
\end{example}

\begin{example}[Foreign key to cross table]
\label{ex:crosstable}
$S_1$ has $A(a, b?)$ with key $\{a\}$, $B(b, c)$ with key $\{b\}$, and the
foreign key $A[b] \subseteq B[b]$. $S_2$ has $A(a)$ with key $\{a\}$, the
cross table $C(a, b)$ with key $\{a\}$, and $B(b, c)$ with key $\{b\}$. The
mapping has two rules:
\begin{align*}
  A(\fv{a}, \fv{b}?) &\corr A(\fv{a}) \wedge \Opt{\fv{b}}(C(\fv{a}, \fv{b})),\\
  B(\fv{b}, \fv{c}) &\corr B(\fv{b}, \fv{c}).
\end{align*}
The absent case of the optional row is $\neg C(\fv{a}, \_)$, and
$\fv{a}$ occurs in two binding atoms on the right, which join on its value. A first
version that joined $B$ into the first rule would leave rows of $B$ that no
$A$ references uncovered on both sides. Accounting finds two constraints
missing in $S_2$: $C[b] \subseteq B[b]$, the image of the foreign key, and
$C[a] \subseteq A[a]$, for $S_2$-covering. The key $\{a\}$ of $C$ is
native: it is the image of the functionality of $b$ in $A$.
\end{example}

\begin{example}[Horizontal partitioning by a discriminator]
\label{ex:horizontal}
$S_1$ has $R(a, b, c)$ with key $\{a\}$ and $c \in \{1, 2\}$. $S_2$ has
$R_1(a, b)$ and $R_2(a, b)$, each with key $\{a\}$, and $R_1[a] \cap
R_2[a] = \emptyset$. The mapping has two rules:
\begin{align*}
  R(\fv{a}, \fv{b}, \dataVal{1}) &\corr R_1(\fv{a}, \fv{b}),\\
  R(\fv{a}, \fv{b}, \dataVal{2}) &\corr R_2(\fv{a}, \fv{b}).
\end{align*}
$S_1$-covering holds because of the enumeration constraint; without it a
row with $c = 3$ would be uncovered. Unfolding the key of $R$ through both
rules gives three cases: both rows in the first rule (the key of $R_1$),
both in the second (the key of $R_2$), and one in each, which gives
$R_1(\fv{a}, \fv{b}) \wedge R_2(\fv{a}, \fv{b}') \to \dataVal{1} \dataEq
\dataVal{2}$, that is the disjointness. The enumeration needs no
counterpart in $S_2$: every row the rules produce has $c \in \{1, 2\}$.
Nothing is missing on either side.
\end{example}

In general, a key whose rows are spread over several rules unfolds into a
key per rule plus a disjointness between the rules.

\section{Open issues}

\begin{enumerate}
  \item Prove Conjecture~\ref{conj:bijection}, or find the extra condition
    it needs (see the note in its proof).
  \item Identification schemes should move into \texttt{agg-logic.tex},
    replacing or generalising $\RefSchemes$.
  \item Characterise a class of constraints (keys, mandatory, subset,
    exclusion over tree patterns) for which key-safety, coverage and
    accounting are decidable, ideally in polynomial time.
  \item Objectified fact types are not covered; in encoding~B they are not
    needed for relational schemas, but ORM2 schemas use them.
  \item Bijective data functions (unit conversion, splitting a name) and
    surrogate identifiers generated on one side only: in the latter case
    the round trip holds up to renaming of the generated values, which
    should be a weaker notion of $\cong$ that also renames those values.
  \item Whether patterns should allow several occurrences of the same
    predicate (self-joins) without restriction; (P1) allows it, but the
    graphical editor needs numbered occurrences.
  \item The relationship with the derivation rules of
    \texttt{agg-logic.tex}: a rule with $\psi$ a single atom and $\varphi$
    without negation reads as a derivation rule for that atom, so
    ``derive, then fold'' can reuse the derivation rule syntax.
\end{enumerate}

\end{document}
